\section{Constitutional Foundations}
\label{sec:constitutional}

\subsection{The Elections Clause}

The constitutional authority for federal regulation of congressional redistricting derives from Article I, Section 4, Clause 1: ``The Times, Places and Manner of holding Elections for Senators and Representatives, shall be prescribed in each State by the Legislature thereof; but the Congress may at any time by Law make or alter such Regulations, except as to the Places of chusing Senators.''\footnote{U.S. Const. art. I, \S\,4, cl. 1.} The text on its face grants Congress plenary authority to ``make or alter'' the regulations states adopt for federal elections, including the manner of conducting them.

The scope of this authority has been broad from the founding. In \textit{Ex parte Siebold} (1879), the Court held that the Clause grants Congress authority to regulate not just procedures but the actual conduct of federal elections, including the appointment and supervision of election officers.\footnote{\textit{Ex parte Siebold}, 100 U.S. 371 (1879).} In \textit{Smiley v.\ Holm} (1932), the Court addressed the congressional redistricting context directly, holding that Congress's power under the Elections Clause includes authority to regulate the drawing of congressional district boundaries---the ``manner'' of elections encompasses the boundaries of the constituencies from which representatives are chosen.\footnote{\textit{Smiley v.\ Holm}, 285 U.S. 355, 366-67 (1932).}

\textit{Smiley} established the crucial precedent: Congress may prescribe not merely \textit{that} states draw districts but \textit{how} they draw them. The Court sustained a federal statute requiring states to establish single-member districts and struck down Minnesota's congressional redistricting statute as inconsistent with federal law. The implication is direct: a federal statute specifying that states must use an algorithmic method to draw congressional districts is an exercise of precisely the power \textit{Smiley} confirmed.

\subsection{One Person, One Vote}

\textit{Wesberry v.\ Sanders} (1964) established that congressional districts must be as equal in population as practicable, deriving this requirement from Article I, Section 2's command that representatives be chosen ``by the People.'' The Court held that ``as nearly as is practicable one man's vote in a congressional election is to be worth as much as another's.''\footnote{\textit{Wesberry v.\ Sanders}, 376 U.S. 1, 7-8 (1964).}

Algorithmic redistricting satisfies this requirement structurally. The algorithm analyzed in \citet{dellaLibera2026recursive} achieves population deviations of less than $\pm 0.5\%$ of the ideal district population for 98\% of districts across all 50 states and three census years (2000, 2010, 2020). This performance meets or exceeds the \textit{Karcher v.\ Daggett} standard for congressional districts, which requires minimizing population deviations unless any remaining deviation is justified by a legitimate state objective.\footnote{\textit{Karcher v.\ Daggett}, 462 U.S. 725 (1983).}

The algorithm achieves this population balance not by relaxing accuracy requirements but by the mathematical structure of recursive bisection: at each level of the bisection tree, the algorithm divides the available population into two groups that differ by at most one unit. Because the target population for each district is derived from the census total divided by the number of districts, the accumulated error across bisection levels is bounded.

\subsection{Post-\textit{Rucho} Federal Role}

\textit{Rucho v.\ Common Cause} is not an obstacle to the federal legislative pathway; it is, paradoxically, its justification. The Court's reasoning explicitly invited congressional action. Chief Justice Roberts wrote: ``Provisions in state statutes and state constitutions can provide standards and guidance for state courts to apply.'' More directly: ``We express no view on any of these state law claims. Nor do we understand our decision to address the authority of state courts to entertain and adjudicate such claims. And courts are not the only check on partisan gerrymandering; other reform mechanisms are available.''\footnote{\textit{Rucho}, 588 U.S. at 721.}

The Court identified Congress as one such mechanism. The dissent went further, noting that ``the majority's reference to other reform mechanisms'' should be read as a signal that ``Congress has the power to act.''\footnote{\textit{Id.} at 748 (Kagan, J., dissenting).} This reading is consistent with the Elections Clause's original purpose: to give Congress a backstop when states fail to maintain the integrity of federal elections. Partisan gerrymandering, which the Court acknowledged is ``unjust'' and ``incompatible with democratic principles,'' is precisely the kind of state failure that the Elections Clause was designed to allow Congress to correct.

\subsection{The Huntington-Hill Precedent}

The most powerful structural precedent for a congressional mandate of algorithmic redistricting is the apportionment statute itself. Since 1941, Congress has mandated that the allocation of House seats among states be computed by the Huntington-Hill equal proportions method, a mathematical formula codified at 2 U.S.C. \S\,2a.\footnote{2 U.S.C. \S\,2a (2022).} Congress enacted this formula to resolve the politically contentious question of how to translate state population counts into House seat allocations---a question that had been decided differently in virtually every decade from 1790 to 1941, with each change benefiting particular states or partisan coalitions.

The analogy to redistricting is structural. Both apportionment and redistricting involve translating population counts into political power. Both invite manipulation when left to political actors---states manipulated apportionment by arguing for alternative mathematical methods that favored their position, just as legislatures manipulate redistricting by drawing favorable boundaries. Both can be governed by mathematical formulas that produce reproducible results from publicly available data. And in both cases, the formula's legitimacy derives not from its optimality on any normative dimension but from its transparency, consistency, and resistance to manipulation.

Congress's choice of the Huntington-Hill method in 1941 was not constitutionally compelled---other methods (Hamilton, Jefferson, Adams, Webster) were mathematically coherent and had historical support. What the choice demonstrated was that Congress has authority to specify, for purposes of federal elections, which mathematical method governs the translation of population counts into political units. A statute mandating the recursive bisection method for congressional districts exercises the same authority in the redistricting context.
